\begin{figure}[!htbp]
\centering
\begin{tikzpicture}[
  font=\sffamily\small,
  module/.style={draw=blue!50!black,rounded corners=3pt,
    text width=2.35cm,minimum height=.67cm,align=center},
  tool/.style={module,text width=1.55cm,minimum height=.65cm,fill=green!14},
  exchange/.style={{Stealth[length=1.6mm]}-{Stealth[length=1.6mm]}},
  request/.style={-{Stealth[length=1.6mm]}}]
\node[module,fill=yellow!18] (core) at (0,0) {智能体核心};
\node[module,fill=blue!12] (plan) at (-3.85,0) {规划};
\node[module,fill=orange!16] (memory) at (3.85,0) {记忆};
\node[module,fill=black!11] (user) at (0,-1.72) {用户};
\node[tool] (web) at (-3.30,1.72) {网页搜索};
\node[tool] (calc) at (-1.10,1.72) {计算器};
\node[tool] (kb) at (1.10,1.72) {知识库};
\node[tool] (other) at (3.30,1.72) {其他工具};
\draw[exchange] (plan) -- (core);
\draw[exchange] (memory) -- (core);
\foreach \toolname in {web,calc,kb,other}{
  \draw[exchange] (core) -- (\toolname.south);
}
\draw[request] ([xshift=-.48cm]user.north) --
  node[left,font=\sffamily\scriptsize]{任务说明} ([xshift=-.48cm]core.south);
\draw[request] ([xshift=.48cm]core.south) --
  node[right,font=\sffamily\scriptsize]{汇总结果} ([xshift=.48cm]user.north);
\node[below,align=left,font=\sffamily\scriptsize] at (plan.south)
  {分解任务\\推理（思维链、Reflexion）};
\node[below,align=left,font=\sffamily\scriptsize] at (memory.south) {短期记忆\\长期记忆};
\draw (-4.17,2.25) -- (-.45,2.25) (0.45,2.25) -- (4.17,2.25);
\draw (-4.17,2.13) -- (-4.17,2.37) (4.17,2.13) -- (4.17,2.37);
\node at (0,2.25) {工具};
\end{tikzpicture}
\caption{大语言模型智能体的主要模块。核心接收用户任务，从记忆模块取得相关信息，交由规划模块分解子任务，再使用可用工具完成子任务。核心最后综合信息并将结果返回用户。}
\label{fig:ch10-agent}
\end{figure}
